\begin{figure}[htbp]
\centering
\begin{tikzpicture}[
  box/.style={draw, rounded corners=2pt, align=center, minimum height=9mm, text width=5.7cm, font=\footnotesize},
  node distance=4mm and 6mm
]
\node[box] (d1) {D01 Everyday life and material culture};
\node[box, right=of d1] (d2) {D02 Language, discourse and pragmatics};
\node[box, below=of d1] (d3) {D03 Social etiquette and interpersonal norms};
\node[box, right=of d3] (d4) {D04 Values, ethics and moral pluralism};
\node[box, below=of d3] (d5) {D05 Law, policy and institutional rules};
\node[box, right=of d5] (d6) {D06 Religion, ritual and taboo};
\node[box, below=of d5] (d7) {D07 Family, kinship, gender and generations};
\node[box, right=of d7] (d8) {D08 Work, education and civic participation};
\node[box, below=of d7] (d9) {D09 Cultural heritage, history, arts and collective memory};
\node[box, right=of d9] (d10) {D10 Identity, diversity and intergroup relations};
\end{tikzpicture}
\caption{The ten equally available diagnostic dimensions. The layout is presentational only and does not impose a hierarchy or fixed parent grouping.}
\label{fig:d01d10-framework}
\end{figure}
